\chapter{Présentation de l'entreprise}
\chapminitoc

\section{Le groupe ALTEN}

L'histoire d'ALTEN commence en 1988, lorsque trois jeunes ingénieurs issus de
grandes écoles, menés par Simon Azoulay --- toujours Président-Directeur Général
du groupe aujourd'hui ---, font un pari : les grands industriels auront de plus
en plus besoin d'externaliser une partie de leur R\&D auprès de partenaires
capables de leur fournir de l'ingénierie de haut niveau. Près de quarante ans
plus tard, ce pari est devenu le \textbf{leader mondial de l'ingénierie et des
services technologiques} (\emph{Engineering and Technology Consulting}).

Les chiffres donnent la mesure du groupe : environ \textbf{4,1 milliards
d'euros} de chiffre d'affaires en 2025, près de \textbf{58\,000 collaborateurs}
--- dont plus de \textbf{51\,000 ingénieurs et consultants} ---, une présence
dans \textbf{plus de 30 pays} et une activité réalisée aux deux tiers à
l'international. Coté à la bourse de Paris (Euronext, indice CAC Mid 60) depuis
1999, ALTEN a bâti cette croissance à la fois organiquement et par une
politique d'acquisitions soutenue en Europe, en Amérique du Nord et en Asie.

% TODO : Figure — carte d'implantation mondiale d'ALTEN (slide dispo sur
% l'intranet ou dans la plaquette « ALTEN Group Overview »)
\begin{figure}[ht]
  \centering
  \fbox{\parbox[c][4.5cm][c]{0.85\textwidth}{\centering\itshape
  Figure à insérer : implantation mondiale du groupe ALTEN\\
  (source : plaquette institutionnelle ALTEN)}}
  \caption{Présence mondiale du groupe ALTEN}
\end{figure}

\section{Secteurs d'activité}

ALTEN accompagne la R\&D des grands acteurs de pratiquement toutes les
industries de pointe : \textbf{aéronautique et spatial}, \textbf{défense et
naval}, \textbf{automobile} --- secteur historique du groupe et cadre de mon
stage ---, ferroviaire, énergie, sciences de la vie, télécommunications,
banque/finance et retail. Cette diversité est une force : les méthodes et
technologies éprouvées dans un secteur (par exemple la sûreté de
fonctionnement aéronautique) irriguent les projets des autres.

Deux grands métiers structurent l'offre : l'\textbf{ingénierie produit et
process} (conception mécanique, électronique, logiciel embarqué, essais) et les
\textbf{IT services} (systèmes d'information, données, cybersécurité). Le
véhicule autonome et connecté, à la croisée exacte de ces deux mondes ---
mécanique, embarqué, réseaux ---, est l'un des axes stratégiques de recherche
du groupe.

\section{Les ALTEN Labs}

Au-delà des projets menés chez ses clients, ALTEN investit dans sa propre
recherche à travers les \textbf{ALTEN Labs} : des structures internes agiles
dont la mission est d'explorer les technologies de rupture, de construire des
démonstrateurs concrets et de faire monter les consultants en compétence sur
les sujets d'avenir (intelligence artificielle, systèmes autonomes, 5G,
jumeaux numériques\ldots). Les Labs fonctionnent comme de petites équipes
projet pluridisciplinaires : un sujet, un encadrant, des stagiaires et
consultants, du matériel, et un objectif de démonstration.

% TODO : reprendre ici 2-3 phrases officielles de la page « Les Labs ALTEN »
% (alten.com) et citer la source en bibliographie.

\section{Le site de Sèvres}

Mon stage s'est déroulé sur le site ALTEN de \textbf{Sèvres (Hauts-de-Seine)},
au 119--121 Grande Rue, à quelques kilomètres du siège du groupe de
Boulogne-Billancourt. Ce site héberge notamment des activités de laboratoire
et les plateformes de démonstration des Labs --- dont les véhicules ProtoVA.

% TODO PHOTOS : prends 2-3 photos sur place (façade du bâtiment, l'atelier /
% la piste d'essai, les deux ProtoVA côte à côte) — bien plus percutant pour
% le jury que n'importe quelle image d'archive.
\begin{figure}[ht]
  \centering
  \fbox{\parbox[c][4.5cm][c]{0.4\textwidth}{\centering\itshape Photo à insérer :
  site ALTEN de Sèvres}}\hfill
  \fbox{\parbox[c][4.5cm][c]{0.4\textwidth}{\centering\itshape Photo à insérer :
  les véhicules ProtoVA}}
  \caption{Le site de Sèvres et la plateforme ProtoVA}
\end{figure}

\section{Le projet ProtoVA}

\textbf{ProtoVA} (\emph{Prototype de Véhicule Autonome}) est le démonstrateur
« véhicule autonome et connecté » des Labs : deux voitures à échelle réduite,
architecturées autour d'un calculateur embarqué NVIDIA Jetson et de ROS~2,
équipées d'un lidar, d'une caméra et d'une chaîne de traction instrumentée.
La plateforme sert de banc d'essai réel --- et non simulé --- aux briques du
véhicule de demain : conduite à distance via 5G, perception et navigation
autonome, communication inter-véhicules selon les normes européennes ITS, et
supervision à distance.

\begin{itemize}
  \item \textbf{Protova1} : première génération (Jetson Nano), utilisée pour
  les développements bas niveau (châssis, motorisation, odométrie) ;
  \item \textbf{Protova2} : seconde génération (Jetson Orin), qui concentre la
  téléopération, l'autonomie et la connectivité.
\end{itemize}

\section{Mon équipe et ma place dans le projet}

% TODO : organigramme (groupe -> direction Labs -> équipe ProtoVA -> toi),
% prénoms/fonctions de ton tuteur et des collègues, comme dans ton rapport 4A.
\begin{figure}[ht]
  \centering
  \fbox{\parbox[c][4cm][c]{0.85\textwidth}{\centering\itshape Figure à insérer :
  organigramme de l'équipe (encadrement ALTEN Labs et stagiaires ProtoVA)}}
  \caption{Organigramme de l'équipe d'accueil}
\end{figure}

J'ai intégré l'équipe ProtoVA en tant que stagiaire ingénieure mécatronique,
avec un périmètre volontairement transverse : mécanique (Protova1), logiciel
embarqué et réseaux (Protova2), et développement web (supervision). Cette
position m'a permis de toucher à toute la chaîne d'un véhicule autonome, du
signal d'encodeur jusqu'au message V2V --- exactement ce que je cherchais dans
ce stage de fin d'études.
